\documentclass[10pt,a4paper]{article}

\usepackage[margin=2.3cm]{geometry}
\usepackage[T1]{fontenc}
\usepackage[utf8]{inputenc}
\usepackage{lmodern}
\usepackage{amsmath,amssymb,amsthm}
\usepackage{tikz-cd}
\usepackage{enumitem}
\usepackage{xcolor}
\usepackage{fancyhdr}

\definecolor{heading}{HTML}{7A1F32}

\pagestyle{fancy}
\fancyhf{}
\lhead{Advanced Software Engineering}
\rhead{Homework 4}
\cfoot{\thepage}
\setlength{\headheight}{14pt}

\newtheorem{lemma}{Lemma}
\newtheorem{theorem}{Theorem}
\theoremstyle{definition}
\newtheorem{definition}{Definition}

\newcommand{\Graph}{\mathbf{Graph}}
\newcommand{\TG}{\mathit{TG}}
\newcommand{\type}{\mathit{type}}
\newcommand{\id}{\mathit{id}}

\title{\textcolor{heading}{\textbf{Homework 4}}}
\author{Matteo Carrese\\Advanced Software Engineering}
\date{}
\begin{document}

\maketitle

\section{Goal}

This homework proves that graphs typed over a fixed type graph $\TG$, together
with type-preserving graph morphisms, form a category. The proof first shows
that graph morphisms are closed under composition. It then verifies
composition, identities, associativity, and the identity laws for typed graphs.


\section{Basic definitions}

\begin{definition}[Directed graph]
A directed graph is a tuple
\[
    G=(V_G,E_G,s_G,t_G),
\]
where $V_G$ is the set of vertices, $E_G$ is the set of edges, and
$s_G,t_G:E_G\to V_G$ give the source and target of every edge.
\end{definition}

\begin{definition}[Graph morphism]
Let $G$ and $H$ be directed graphs. A graph morphism $f:G\to H$ is a pair
of functions
\[
    f_V:V_G\to V_H,
    \qquad
    f_E:E_G\to E_H
\]
that preserve the source and target of every edge:
\begin{equation}
    f_V\circ s_G=s_H\circ f_E,
    \qquad
    f_V\circ t_G=t_H\circ f_E.
    \label{eq:graph-morphism}
\end{equation}
\end{definition}

The two equations say that mapping an edge and then reading its endpoint gives
the same result as reading the original endpoint and then mapping that vertex.
Therefore, a graph morphism preserves the incidence structure of the graph.

\section{Composition of graph morphisms}

The following lemma proves the explicit homework statement appearing earlier
in the lecture material.

\begin{lemma}
The composition of two graph morphisms is a graph morphism.
\end{lemma}

\begin{proof}
Let
\[
    f=(f_V,f_E):G\to H
    \qquad\text{and}\qquad
    g=(g_V,g_E):H\to K
\]
be graph morphisms. Define their composition componentwise:
\[
    g\circ f=(g_V\circ f_V,\;g_E\circ f_E):G\to K.
\]

We first verify preservation of sources. Using associativity of function
composition and the morphism equations for $f$ and $g$, we obtain
\begin{align*}
 (g_V\circ f_V)\circ s_G
 &=g_V\circ(f_V\circ s_G)\\
 &=g_V\circ(s_H\circ f_E)\\
 &=(g_V\circ s_H)\circ f_E\\
 &=(s_K\circ g_E)\circ f_E\\
 &=s_K\circ(g_E\circ f_E).
\end{align*}

The same calculation for targets gives
\[
    (g_V\circ f_V)\circ t_G
    =t_K\circ(g_E\circ f_E).
\]
Hence $g\circ f$ satisfies both equations in
Equation~\ref{eq:graph-morphism}, so it is a graph morphism.
\end{proof}

\section{Typed graphs}

Fix one graph $\TG$, called the \emph{type graph}.

\begin{definition}[Typed graph]
A graph typed over $\TG$ is a pair $(G,\type_G)$, where $G$ is a graph and
\[
    \type_G:G\to\TG
\]
is a graph morphism. It assigns a type in $\TG$ to every vertex and edge of
$G$ while preserving sources and targets.
\end{definition}

\begin{definition}[Type-preserving morphism]
A morphism between two graphs typed over the same $\TG$,
\[
    f:(G,\type_G)\to(H,\type_H),
\]
is a graph morphism $f:G\to H$ satisfying
\begin{equation}
    \type_H\circ f=\type_G.
    \label{eq:type-preservation}
\end{equation}
This equality means that the following triangle commutes:
\[
\begin{tikzcd}[column sep=large,row sep=large]
G \arrow[rr,"f"] \arrow[dr,swap,"\type_G"]
  & & H \arrow[dl,"\type_H"] \\
  & \TG &
\end{tikzcd}
\]
\end{definition}

In simple terms, $f$ may map vertices and edges to different elements, but it
cannot change their types.

\section{The category of graphs typed over $\TG$}

We define a structure denoted by $\Graph/\TG$:

\begin{itemize}[leftmargin=2em]
    \item its \textbf{objects} are all pairs $(G,\type_G)$ with
          $\type_G:G\to\TG$;
    \item its \textbf{morphisms} are graph morphisms that satisfy
          Equation~\ref{eq:type-preservation};
    \item composition is the componentwise composition of the underlying
          graph morphisms;
    \item the identity on $(G,\type_G)$ is the ordinary graph identity
          $\id_G$.
\end{itemize}

\begin{theorem}
The structure $\Graph/\TG$ is a category.
\end{theorem}

\begin{proof}
We verify the required properties.

\paragraph{Closure under composition.}
Consider two composable type-preserving morphisms
\[
 f:(G,\type_G)\to(H,\type_H),
 \qquad
 g:(H,\type_H)\to(K,\type_K).
\]
By the lemma, the underlying composition $g\circ f:G\to K$ is a graph
morphism. It also preserves typing because
\begin{align*}
 \type_K\circ(g\circ f)
 &=(\type_K\circ g)\circ f
 &&\text{by associativity}\\
 &=\type_H\circ f
 &&\text{because $g$ preserves types}\\
 &=\type_G
 &&\text{because $f$ preserves types}.
\end{align*}
Therefore, $g\circ f$ is a morphism in $\Graph/\TG$.

The corresponding commuting diagram is
\[
\begin{tikzcd}[column sep=large,row sep=large]
G \arrow[r,"f"] \arrow[rr,bend left=35,"g\circ f"]
  \arrow[dr,swap,"\type_G"]
& H \arrow[r,"g"] \arrow[d,"\type_H"]
& K \arrow[dl,"\type_K"] \\
& \TG &
\end{tikzcd}
\]

\paragraph{Identity morphisms.}
For every object $(G,\type_G)$, the graph identity
$\id_G:G\to G$ preserves typing:
\[
    \type_G\circ\id_G=\type_G.
\]
Thus, every object has a valid identity morphism in $\Graph/\TG$.

\paragraph{Associativity.}
Typed morphisms are composed through their vertex and edge functions. Function
composition is associative, so for all composable typed morphisms $f$, $g$,
and $h$,
\[
    (h\circ g)\circ f=h\circ(g\circ f).
\]
Closure, proved above, ensures that both sides are still type-preserving
morphisms.

\paragraph{Identity laws.}
For every typed morphism
$f:(G,\type_G)\to(H,\type_H)$, the identity laws of the underlying graph
morphisms give
\[
    \id_H\circ f=f
    \qquad\text{and}\qquad
    f\circ\id_G=f.
\]
The identities and $f$ preserve typing, so these equalities hold inside
$\Graph/\TG$.

All the category requirements are satisfied. Therefore, graphs typed over the
fixed graph $\TG$, with type-preserving graph morphisms, form a category.
\end{proof}

\section{Categorical interpretation}

The category just constructed is the \emph{slice category} of graphs over
$\TG$. An object is an arrow $G\to\TG$, and a morphism is an arrow between
their domains that makes the typing triangle commute. This is commonly written
as $\Graph/\TG$ and is also represented in the lecture slides as
$\Graph\downarrow\TG$.

This interpretation explains the construction at a more general level: for
any category $\mathbf{C}$ and fixed object $X$, the arrows with codomain $X$
and their commuting triangles form the slice category $\mathbf{C}/X$.

\section{Conclusion}

The essential point is that all operations remain compatible with the fixed
typing graph. Ordinary graph morphisms compose, and the equation
$\type_H\circ f=\type_G$ remains true after composition. Identity morphisms
also preserve typing, while associativity and the identity laws are inherited
componentwise from ordinary functions. Consequently, $\Graph/\TG$ is a
well-defined category.

\end{document}\documentclass[10pt,a4paper]{article}

\usepackage[margin=2.3cm]{geometry}
\usepackage[T1]{fontenc}
\usepackage[utf8]{inputenc}
\usepackage{lmodern}
\usepackage{amsmath,amssymb,amsthm}
\usepackage{tikz-cd}
\usepackage{enumitem}
\usepackage{xcolor}
\usepackage{fancyhdr}

\definecolor{heading}{HTML}{7A1F32}

\pagestyle{fancy}
\fancyhf{}
\lhead{Advanced Software Engineering}
\rhead{Homework 4}
\cfoot{\thepage}
\setlength{\headheight}{14pt}

\newtheorem{lemma}{Lemma}
\newtheorem{theorem}{Theorem}
\theoremstyle{definition}
\newtheorem{definition}{Definition}

\newcommand{\Graph}{\mathbf{Graph}}
\newcommand{\TG}{\mathit{TG}}
\newcommand{\type}{\mathit{type}}
\newcommand{\id}{\mathit{id}}

\title{\textcolor{heading}{\textbf{Homework 4}}}
\author{Matteo Carrese\\Advanced Software Engineering}
\date{}
\begin{document}

\maketitle

\section{Goal}

This homework proves that graphs typed over a fixed type graph $\TG$, together
with type-preserving graph morphisms, form a category. The proof first shows
that graph morphisms are closed under composition. It then verifies
composition, identities, associativity, and the identity laws for typed graphs.


\section{Basic definitions}

\begin{definition}[Directed graph]
A directed graph is a tuple
\[
    G=(V_G,E_G,s_G,t_G),
\]
where $V_G$ is the set of vertices, $E_G$ is the set of edges, and
$s_G,t_G:E_G\to V_G$ give the source and target of every edge.
\end{definition}

\begin{definition}[Graph morphism]
Let $G$ and $H$ be directed graphs. A graph morphism $f:G\to H$ is a pair
of functions
\[
    f_V:V_G\to V_H,
    \qquad
    f_E:E_G\to E_H
\]
that preserve the source and target of every edge:
\begin{equation}
    f_V\circ s_G=s_H\circ f_E,
    \qquad
    f_V\circ t_G=t_H\circ f_E.
    \label{eq:graph-morphism}
\end{equation}
\end{definition}

The two equations say that mapping an edge and then reading its endpoint gives
the same result as reading the original endpoint and then mapping that vertex.
Therefore, a graph morphism preserves the incidence structure of the graph.

\section{Composition of graph morphisms}

The following lemma proves the explicit homework statement appearing earlier
in the lecture material.

\begin{lemma}
The composition of two graph morphisms is a graph morphism.
\end{lemma}

\begin{proof}
Let
\[
    f=(f_V,f_E):G\to H
    \qquad\text{and}\qquad
    g=(g_V,g_E):H\to K
\]
be graph morphisms. Define their composition componentwise:
\[
    g\circ f=(g_V\circ f_V,\;g_E\circ f_E):G\to K.
\]

We first verify preservation of sources. Using associativity of function
composition and the morphism equations for $f$ and $g$, we obtain
\begin{align*}
 (g_V\circ f_V)\circ s_G
 &=g_V\circ(f_V\circ s_G)\\
 &=g_V\circ(s_H\circ f_E)\\
 &=(g_V\circ s_H)\circ f_E\\
 &=(s_K\circ g_E)\circ f_E\\
 &=s_K\circ(g_E\circ f_E).
\end{align*}

The same calculation for targets gives
\[
    (g_V\circ f_V)\circ t_G
    =t_K\circ(g_E\circ f_E).
\]
Hence $g\circ f$ satisfies both equations in
Equation~\ref{eq:graph-morphism}, so it is a graph morphism.
\end{proof}

\section{Typed graphs}

Fix one graph $\TG$, called the \emph{type graph}.

\begin{definition}[Typed graph]
A graph typed over $\TG$ is a pair $(G,\type_G)$, where $G$ is a graph and
\[
    \type_G:G\to\TG
\]
is a graph morphism. It assigns a type in $\TG$ to every vertex and edge of
$G$ while preserving sources and targets.
\end{definition}

\begin{definition}[Type-preserving morphism]
A morphism between two graphs typed over the same $\TG$,
\[
    f:(G,\type_G)\to(H,\type_H),
\]
is a graph morphism $f:G\to H$ satisfying
\begin{equation}
    \type_H\circ f=\type_G.
    \label{eq:type-preservation}
\end{equation}
This equality means that the following triangle commutes:
\[
\begin{tikzcd}[column sep=large,row sep=large]
G \arrow[rr,"f"] \arrow[dr,swap,"\type_G"]
  & & H \arrow[dl,"\type_H"] \\
  & \TG &
\end{tikzcd}
\]
\end{definition}

In simple terms, $f$ may map vertices and edges to different elements, but it
cannot change their types.

\section{The category of graphs typed over $\TG$}

We define a structure denoted by $\Graph/\TG$:

\begin{itemize}[leftmargin=2em]
    \item its \textbf{objects} are all pairs $(G,\type_G)$ with
          $\type_G:G\to\TG$;
    \item its \textbf{morphisms} are graph morphisms that satisfy
          Equation~\ref{eq:type-preservation};
    \item composition is the componentwise composition of the underlying
          graph morphisms;
    \item the identity on $(G,\type_G)$ is the ordinary graph identity
          $\id_G$.
\end{itemize}

\begin{theorem}
The structure $\Graph/\TG$ is a category.
\end{theorem}

\begin{proof}
We verify the required properties.

\paragraph{Closure under composition.}
Consider two composable type-preserving morphisms
\[
 f:(G,\type_G)\to(H,\type_H),
 \qquad
 g:(H,\type_H)\to(K,\type_K).
\]
By the lemma, the underlying composition $g\circ f:G\to K$ is a graph
morphism. It also preserves typing because
\begin{align*}
 \type_K\circ(g\circ f)
 &=(\type_K\circ g)\circ f
 &&\text{by associativity}\\
 &=\type_H\circ f
 &&\text{because $g$ preserves types}\\
 &=\type_G
 &&\text{because $f$ preserves types}.
\end{align*}
Therefore, $g\circ f$ is a morphism in $\Graph/\TG$.

The corresponding commuting diagram is
\[
\begin{tikzcd}[column sep=large,row sep=large]
G \arrow[r,"f"] \arrow[rr,bend left=35,"g\circ f"]
  \arrow[dr,swap,"\type_G"]
& H \arrow[r,"g"] \arrow[d,"\type_H"]
& K \arrow[dl,"\type_K"] \\
& \TG &
\end{tikzcd}
\]

\paragraph{Identity morphisms.}
For every object $(G,\type_G)$, the graph identity
$\id_G:G\to G$ preserves typing:
\[
    \type_G\circ\id_G=\type_G.
\]
Thus, every object has a valid identity morphism in $\Graph/\TG$.

\paragraph{Associativity.}
Typed morphisms are composed through their vertex and edge functions. Function
composition is associative, so for all composable typed morphisms $f$, $g$,
and $h$,
\[
    (h\circ g)\circ f=h\circ(g\circ f).
\]
Closure, proved above, ensures that both sides are still type-preserving
morphisms.

\paragraph{Identity laws.}
For every typed morphism
$f:(G,\type_G)\to(H,\type_H)$, the identity laws of the underlying graph
morphisms give
\[
    \id_H\circ f=f
    \qquad\text{and}\qquad
    f\circ\id_G=f.
\]
The identities and $f$ preserve typing, so these equalities hold inside
$\Graph/\TG$.

All the category requirements are satisfied. Therefore, graphs typed over the
fixed graph $\TG$, with type-preserving graph morphisms, form a category.
\end{proof}

\section{Categorical interpretation}

The category just constructed is the \emph{slice category} of graphs over
$\TG$. An object is an arrow $G\to\TG$, and a morphism is an arrow between
their domains that makes the typing triangle commute. This is commonly written
as $\Graph/\TG$ and is also represented in the lecture slides as
$\Graph\downarrow\TG$.

This interpretation explains the construction at a more general level: for
any category $\mathbf{C}$ and fixed object $X$, the arrows with codomain $X$
and their commuting triangles form the slice category $\mathbf{C}/X$.

\section{Conclusion}

The essential point is that all operations remain compatible with the fixed
typing graph. Ordinary graph morphisms compose, and the equation
$\type_H\circ f=\type_G$ remains true after composition. Identity morphisms
also preserve typing, while associativity and the identity laws are inherited
componentwise from ordinary functions. Consequently, $\Graph/\TG$ is a
well-defined category.

\end{document}

